%% =====================================================================
%% BLG 475E Software Quality and Testing - Term Project
%% Phase 1 Report - IEEE Journal Template (per assignment brief)
%% Authors: Mustafa Eren KOC (150190805), Onat Baris Ercan (150210075)
%% =====================================================================
\documentclass[journal]{IEEEtran}

\usepackage{amsmath}
\usepackage{amssymb}
\usepackage{graphicx}
\usepackage{booktabs}
\usepackage{multirow}
\usepackage{array}
\usepackage{url}
\usepackage{listings}
\usepackage{xcolor}
\usepackage{hyperref}

\hypersetup{
    colorlinks=true,
    linkcolor=blue,
    citecolor=blue,
    urlcolor=cyan,
    pdftitle={BLG 475E Term Project - Phase 1 Report}
}

% --- Code listing style (compact for two-column IEEE layout) ---
\definecolor{codegreen}{rgb}{0,0.6,0}
\definecolor{codegray}{rgb}{0.5,0.5,0.5}
\definecolor{codepurple}{rgb}{0.58,0,0.82}
\definecolor{backcolour}{rgb}{0.96,0.96,0.92}

\lstdefinestyle{ieeeJava}{
    backgroundcolor=\color{backcolour},
    commentstyle=\color{codegreen},
    keywordstyle=\color{magenta},
    numberstyle=\tiny\color{codegray},
    stringstyle=\color{codepurple},
    basicstyle=\ttfamily\scriptsize,
    breakatwhitespace=false,
    breaklines=true,
    captionpos=b,
    keepspaces=true,
    numbers=left,
    numbersep=5pt,
    showspaces=false,
    showstringspaces=false,
    showtabs=false,
    tabsize=2,
    columns=fullflexible,
    xleftmargin=10pt,
    aboveskip=4pt,
    belowskip=4pt
}
\lstset{style=ieeeJava}

\hyphenation{op-tical net-works semi-conduc-tor Hu-man-Eval Java-Script Power-Shell}

\begin{document}

\title{Evaluating Local Code-LLMs on HumanEval-X Java Prompts:\\
Semi-Agentic Generation, Improved Tests, and Mutation-Based Guardrails}

\author{Mustafa~Eren~Ko\c{c}~(150190805),~and~Onat~Bar{\i}\c{s}~Ercan~(150210075)% <-this % stops a space
\thanks{The authors are with the Department of Computer Engineering,
Istanbul Technical University, Istanbul, T\"urkiye, e-mail:
\{koc20, ercano21\}@itu.edu.tr.}% <-this % stops a space
\thanks{This report is prepared for the BLG~475E Software Quality and
Testing course (2025--2026 Spring term), supervised by Dr.~Ay\c{s}e~Tosun~K\"uhn.
The authors form a two-person group with the explicit acknowledgement of the
course instructor. Source code, scripts, interaction logs and CI artefacts
are publicly available at \url{https://github.com/mustafaerenkoc44/BLG-475E-Term-Project}.}}

\markboth{BLG~475E Software Quality and Testing, Phase~1 Report, Spring 2026}%
{Ko\c{c} and Ercan: Evaluating Local Code-LLMs on HumanEval-X Java Prompts}

\maketitle

\begin{abstract}
This paper reports on Phase~1 of the BLG~475E term project, which studies how
two publicly available local code-oriented Large Language Models (LLMs),
\texttt{Qwen2.5-Coder-1.5B-Instruct} and \texttt{DeepSeek-Coder-1.3B-Instruct},
both executed through the same \texttt{llama.cpp} runtime, perform on a
30-prompt subset of the HumanEval-X Java corpus. Each prompt is processed
through a \emph{semi-agentic} pipeline in which prompt submission, raw
output, text normalization, dataset base tests, JUnit~6 improved tests,
JaCoCo branch coverage and hand-crafted mutation-based tests are all
recorded as reproducible artefacts. After at most one logged refactoring
round both models reach $30/30$ base-test success; the improved suites lift
aggregate branch coverage from $96.09\,\%$ / $98.46\,\%$ (Qwen / DeepSeek)
on the base suites to $98.44\,\%$ / $100.00\,\%$. Every (task, model) pair
in the corpus carries at least one \texttt{improvedMutation\ldots} test
that encodes a specific operator mutation (ROR, AOR, NPE, BND, RV, LCR,
UOI, SBR or EXC), giving a per-suite mutation score of $100\,\%$ for the
improved suites and $0\,\%$ for the LLM-produced base suites. We argue
that branch coverage alone is insufficient to evaluate LLM test generators
and that a structured equivalence-partitioning plus mutation-kill trace
materially improves the transparency of the benchmark. The methodology
established in this phase carries directly into the Phase~2 extension that
lifts the selected helpers (\texttt{Java/18}, \texttt{Java/23},
\texttt{Java/27}) into a composite \texttt{BookScan} class.
\end{abstract}

\begin{IEEEkeywords}
Large language models, unit testing, JUnit, JaCoCo, mutation testing,
equivalence partitioning, HumanEval-X, semi-agentic workflow, software
quality.
\end{IEEEkeywords}

% =====================================================================
\section{Introduction}\label{sec:intro}
% =====================================================================
\IEEEPARstart{M}{odern} software teams increasingly use Large Language
Models (LLMs) to draft both source code and its
tests~\cite{chatunitest,metatestgen,lit_review}. The economic appeal is
obvious --- a model can produce compilable code and candidate tests in
seconds --- but the engineering risk is equally clear: generated code can
compile and pass a happy-path test without actually meeting the
specification, and generated tests can execute without exercising any
meaningful behaviour. BLG~475E \emph{Software Quality and Testing} asks
students to investigate both sides of that trade-off on a controlled Java
corpus. Phase~1 of the project, reported in this paper, establishes the
methodology, the toolchain and the baseline measurements; Phase~2,
reported separately, builds on these foundations to study integration
testing of a composite class.

The mandate of Phase~1, taken from the course brief, is to:
\begin{enumerate}
  \item select 30 HumanEval-X Java prompts balanced across easy /
    moderate / hard categories;
  \item generate both code and an initial test suite with two
    open-weight code-LLMs under identical prompting conditions;
  \item run the dataset base tests and extract JaCoCo branch coverage;
  \item \emph{improve} the tests using equivalence-class reasoning,
    boundary-value analysis, and mutation-style guards;
  \item log every LLM interaction and document repairs;
  \item review the recent literature and produce an analysis report.
\end{enumerate}

The contribution of this paper is twofold. First, we show that an orthodox
semi-agentic pipeline --- PowerShell + Python drivers, \texttt{llama.cpp}
inference, JUnit~6 runtime, JaCoCo CLI --- is sufficient to reproduce every
experiment on a student laptop and on a GitHub Actions runner. Second, we
add a hand-crafted \emph{mutation-kill matrix} that complements branch
coverage and explains why two LLMs with similar coverage numbers can still
have visibly different robustness profiles, e.g.\ Qwen's \texttt{Math.floor}
versus DeepSeek's \texttt{(int)} cast for negative truncation
(Section~\ref{sec:results}).

The remainder of the paper is organised as follows. Section~\ref{sec:method}
details the methodology; Section~\ref{sec:challenges} documents the technical
challenges encountered (Unicode workspace paths, llama.cpp boilerplate echo,
repair loops); Section~\ref{sec:results} presents the measured results;
Section~\ref{sec:discuss} discusses their implications and threats to
validity; Section~\ref{sec:lit} contains the five-paper literature review;
Section~\ref{sec:conclusion} concludes Phase~1.

% =====================================================================
\section{Methodology}\label{sec:method}
% =====================================================================

\subsection{LLM Selection}

We adopted a \emph{semi-agentic} workflow because every transition between
steps remains observable, which makes the report defensible and the
experiment reproducible: each step writes its inputs and outputs into the
repository before the next step starts. A fully autonomous loop is not
necessary because the local \texttt{llama.cpp} runtime would have required
additional infrastructure unrelated to the research question. We discuss
this trade-off again in Section~\ref{sec:discuss}.

The two selected models satisfy five constraints (public release;
code-oriented fine-tune; locally deployable in $<$8\,GB; instruction-tuned;
distinct lineage):
\begin{itemize}
    \item \texttt{Qwen/Qwen2.5-Coder-1.5B-Instruct} (HuggingFace), executed
        through the GGUF deployment \texttt{Qwen2.5-Coder-1.5B-Instruct-GGUF}.
    \item \texttt{deepseek-ai/deepseek-coder-1.3b-instruct} (HuggingFace),
        executed through the GGUF deployment \texttt{DeepSeek-Coder-1.3B-Instruct-GGUF}.
\end{itemize}
Both models are loaded into the same local \texttt{llama.cpp} runtime with
identical sampling parameters so that the only varying factor between two
paired experiments is the model itself.

\subsection{Workflow}

For every selected prompt, the pipeline executes nine steps in order:
prompt issue, raw response capture, normalization, base-test execution,
optional repair loop, JUnit conversion, base coverage measurement, test
improvement, and improved coverage measurement. Each step writes a
versioned artefact into \texttt{Phase1/}; GitHub Actions
(\texttt{phase1-ci.yml}) re-runs the entire pipeline on every push.

\subsection{Prompt Selection}

The HumanEval-X Java corpus contains 164 prompts. From these we selected 30
recorded in \texttt{Phase1/data/selected\_prompts.csv}, balanced as 10 easy
/ 10 moderate / 10 hard (Table~\ref{tab:difficulty}). The criteria favoured
domain spread (numeric, string-filtering, parsing, number-theory, list
manipulation, state tracking), boundary sensitivity, and continuity with
Phase~2 (the three tasks Java/18, Java/23 and Java/27 are explicitly
retained because Phase~2 lifts their helpers into a composite
\texttt{BookScan} class). The complete task list is reproduced in
Appendix~\ref{app:prompts}.

\begin{table}[!t]
\renewcommand{\arraystretch}{1.2}
\caption{Phase 1 prompt difficulty distribution.}
\label{tab:difficulty}
\centering
\begin{tabular}{@{}lc@{}}
\toprule
\textbf{Difficulty} & \textbf{Count} \\
\midrule
easy     & 10 \\
moderate & 10 \\
hard     & 10 \\
\midrule
\textbf{total} & \textbf{30} \\
\bottomrule
\end{tabular}
\end{table}

\subsection{Code Generation}

For each prompt the verbatim HumanEval-X declaration is sent to the model.
The course brief forbids editing the prompt at the generation stage, so the
only post-processing is structural normalization: removing Markdown fences,
stripping a re-stated prompt that some models echo, and dropping malformed
fragments that would prevent compilation. Semantic normalization is
forbidden at this step; if a normalized solution still fails the base
test, that failure is recorded and the repair happens in the refactoring
loop, not silently in the normalizer. The full prompt and the full raw
response are stored in \texttt{logs/<model>/Java\_<N>\_code-generation.md}
for every (task, model) pair, giving 60 verbatim code-generation logs.

\begin{lstlisting}[language=Java, caption=Required header block prepended to every generated solution and test file.]
/* @Authors
 * Student Names: Mustafa Eren KOC, Onat Baris ERCAN
 * Student IDs: 150190805, 150210075
 */
\end{lstlisting}

\subsection{Base Test Conversion}

The HumanEval-X reference test ships as a single \texttt{Main.java} driver
that asserts a list of expected return values. To attach a coverage agent
and to integrate cleanly with our reporting tools, every reference test is
re-encoded as a JUnit~6 test class \texttt{DatasetBaseTest.java}. The
conversion is purely structural: each \texttt{assertEquals}-style line in
the driver becomes a single \texttt{@Test} method whose body asserts the
same expected value, and only the syntactic tweaks the brief explicitly
permits at this step are applied. No oracle is changed.

\subsection{Test Improvement}\label{sec:improved}

For every (task, model) pair an improved JUnit suite
\texttt{DatasetImprovedTest.java} is authored on top of the base test.
Three orthogonal sources drive the improvement.

\subsubsection{Test Smell Audit}

Each LLM-generated base suite is hand-audited against the Phase~1
test-smell checklist \cite{jnose}. The most frequent smells are
\emph{happy-path only}, \emph{conditional logic in tests}, \emph{missing
boundary cases}, and \emph{assertion roulette}. The improved suites
explicitly remove these smells.

\subsubsection{Equivalence Partitioning and Boundary-Value Analysis}

For every task we wrote an \texttt{equivalence\_classes.md} document with
three tables: valid classes, invalid classes and boundary conditions. Each
row is annotated with whether the existing test already covers the class.
Rows marked \emph{No} or \emph{Partially} drive at least one extra
\texttt{@Test} method in the improved suite. Table~\ref{tab:eq-example}
reproduces the equivalence-class table for Java/0
(\texttt{hasCloseElements}) as a representative example.

\begin{table}[!t]
\renewcommand{\arraystretch}{1.15}
\caption{Equivalence partitioning for Java/0 (\texttt{hasCloseElements}).
V$=$valid, I$=$invalid, B$=$boundary.}
\label{tab:eq-example}
\centering
\footnotesize
\begin{tabular}{@{}lp{3.0cm}p{2.4cm}c@{}}
\toprule
\textbf{ID} & \textbf{Description} & \textbf{Input} & \textbf{Cov.} \\
\midrule
V1 & At least one pair closer than $\tau$ & \texttt{[1,2.8,3,4]}, $0.3$ & yes \\
V2 & No pair closer than $\tau$           & \texttt{[1,2,3]}, $0.5$     & yes \\
V3 & Duplicates at zero distance          & \texttt{[2,2,5]}, $0.1$     & impr. \\
I1 & Null list                            & \texttt{null}, $0.3$        & impr. \\
I2 & Negative threshold                   & \texttt{[1,2]}, $-0.1$      & impr. \\
B1 & Size below two                       & \texttt{[]} or \texttt{[1]} & impr. \\
B2 & Distance equals $\tau$               & \texttt{[1,1.3]}, $0.3$     & impr. \\
\bottomrule
\end{tabular}
\end{table}

\subsection{Hand-Crafted Mutation Guards}\label{sec:mutation}

Phase~1 explicitly disallows external mutation tooling such as PITest. We
therefore hand-craft mutants and encode them as JUnit assertions: every
improved suite has at least one method whose name starts with
\texttt{improvedMutation\ldots}, and that method is designed so that it
fails if a specific operator-level mutation were introduced into the
corresponding \texttt{Solution.java}. Table~\ref{tab:mutation-catalogue}
summarises the operator catalogue used in Phase~1; the per-task mapping is
maintained in \texttt{Phase1/docs/analysis/mutation\_testing\_strategy.md}.

\begin{table}[!t]
\renewcommand{\arraystretch}{1.15}
\caption{Hand-crafted mutation operators exercised in Phase~1.}
\label{tab:mutation-catalogue}
\centering
\footnotesize
\begin{tabular}{@{}lp{1.9cm}p{3.5cm}@{}}
\toprule
\textbf{Op} & \textbf{Family} & \textbf{Representative use} \\
\midrule
ROR & Relational op.        & Java/3 \verb|< 0| vs \verb|<= 0| \\
AOR & Arithmetic op.        & Java/2 \texttt{Math.floor} vs \texttt{(int)} \\
LCR & Logical connector     & Java/9 \verb|null  ||  size==0| \\
UOI & Unary insertion       & Java/31 negative-prime guard \\
BND & Boundary literal      & Java/25 / 36 numeric guards \\
SBR & Statement removal     & Java/0 self-pair skip \\
RV  & Return-value swap     & Java/39 non-positive index \\
NPE & Null replacement      & every I1 invalid class \\
EXC & Exception flip        & Java/11 unequal-length inputs \\
\bottomrule
\end{tabular}
\end{table}

\begin{lstlisting}[language=Java, caption=Hand-crafted mutation guard for Java/0 (excerpt).]
@Test
void improvedMutationNegativeValuesAndNegativeThreshold() {
    Solution s = new Solution();
    Assertions.assertTrue(
        s.hasCloseElements(
            new ArrayList<>(Arrays.asList(-1.0,-1.05,3.0)), 0.1),
        "Mutation: negative values must preserve "
        + "the absolute-distance oracle");
    Assertions.assertFalse(
        s.hasCloseElements(
            new ArrayList<>(Arrays.asList(1.0,2.0)), -0.1),
        "I2: a negative threshold must reject every pair");
}
\end{lstlisting}

\subsection{Refactoring Loop}\label{sec:refactor}

If a normalized solution fails the dataset base test, the workflow does not
silently fix the code. It issues a refactoring prompt to the same model,
attaches the failure log and the broken \texttt{Solution.java}, and stores
the resulting patch as the new solution. Seven tasks (Java/10, /12, /17,
/19, /23, /27, /39) needed exactly one refactoring round before reaching
\texttt{compile\_success=true} and \texttt{junit\_success=true}; all seven
were eventually fixed and now compile and pass on both models. Every
refactoring round is recorded as a Markdown log under
\texttt{Phase1/logs/<model>/Java\_<N>\_refactoring.md}.

\subsection{Branch Coverage Pipeline}

Branch coverage is measured with JaCoCo~0.8.12 attached to the JUnit~6
console runner (\texttt{junit-platform-console-standalone-6.0.3.jar}). The
Python driver \texttt{run\_base\_coverage.py} compiles each
\texttt{Solution.java} together with the corresponding test class, runs the
suite under the JaCoCo agent, renders the resulting \texttt{jacoco.exec}
into a CSV via the JaCoCo CLI, and appends one row per task to a master
CSV. Two master CSVs are produced for the LLM-generated base suite
(\texttt{Phase1/results/base\_coverage\_results.csv}) and for the improved
suite (\texttt{Phase1/results/improved\_coverage\_results.csv}).

% =====================================================================
\section{Technical Challenges and Solutions}\label{sec:challenges}
% =====================================================================

\subsection{Unicode Workspace Path Issues}

The development workspace lives under a path containing the Turkish
character~\emph{\.{I}} (\verb|D:\Indirilenler\BLG_475E_Term_Project\|).
Several tools mishandled the non-ASCII path: the JaCoCo agent silently
zeroed-out per-task counts when launched with an absolute
\texttt{destfile=} pointer, and the Maven local repository defaulted to a
path that intermittently broke the resolver.

\subsubsection{Solution}
The pipeline now uses a relative \texttt{jacoco.exec} next to the per-task
work directory, redirects the Maven JaCoCo plugin's \texttt{destFile} to
the system tmpdir, switches the base-test result CSV from absolute to
repo-relative paths, and pins the Maven local repository to a
workspace-local \texttt{.m2} folder via \texttt{-Dmaven.repo.local}. After
these fixes the pipeline runs identically under both ASCII and non-ASCII
workspaces.

\subsection{LLM Response Boilerplate Echo}

The local \texttt{llama.cpp} runs for the test-improvement step were
explicitly instructed to \emph{not} restate the JUnit boilerplate they had
been given as context. Both models restated that boilerplate repeatedly
anyway. Reproducing the verbatim raw response in the Markdown log would
have added hundreds of lines of duplicated boilerplate.

\subsubsection{Solution}
The 60 code-generation logs and the 14 refactoring logs continue to record
the verbatim raw response. The 60 test-improvement logs record an
\emph{essential captured shape}: the full prompt verbatim, a summary of
the model's distinct suggestions, a triage table classifying every
suggestion as kept, rewritten or rejected, the final additions that landed
in \texttt{DatasetImprovedTest.java}, and a sanity-check pointer to the CI
row that re-executes the suite. The capture policy is documented at the
top of \texttt{Phase1/logs/README.md}.

\subsection{Repair Loop for Failing Generations}

Seven of the 30 selected tasks initially failed their base test on at
least one model. Table~\ref{tab:repair} summarises them.

\begin{table}[!t]
\renewcommand{\arraystretch}{1.15}
\caption{Tasks that required a refactoring round before reaching base-test success.}
\label{tab:repair}
\centering
\footnotesize
\begin{tabular}{@{}lp{6.0cm}@{}}
\toprule
\textbf{Task} & \textbf{Initial defect} \\
\midrule
Java/10 & \texttt{makePalindrome} returned \texttt{null} when no
        palindromic suffix was found instead of falling back to the input. \\
Java/12 & \texttt{longest} did not handle the empty-list case. \\
Java/17 & \texttt{parseMusic} mis-tokenised inputs with adjacent shorthand
        symbols. \\
Java/19 & \texttt{sortNumbers} produced incorrect ordering for inputs
        containing \texttt{"zero"}. \\
Java/23 & \texttt{strlen} produced an off-by-one error on the Qwen variant. \\
Java/27 & \texttt{flipCase} dropped non-letter characters instead of
        preserving them. \\
Java/39 & \texttt{primeFib} timed out on the larger fixed input under the
        default settings. \\
\bottomrule
\end{tabular}
\end{table}

% =====================================================================
\section{Results}\label{sec:results}
% =====================================================================

All numbers in this section can be reproduced by running the Phase~1
GitHub Actions workflow on a clean checkout; the headline CSVs are
committed under \texttt{Phase1/results/}.

\subsection{Code Generation Outcomes}

After the refactoring loop, both models compile and pass their base tests
on all 30 prompts (Table~\ref{tab:codegen}).

\begin{table}[!t]
\renewcommand{\arraystretch}{1.15}
\caption{Phase 1 code generation outcomes after the refactoring loop.}
\label{tab:codegen}
\centering
\footnotesize
\begin{tabular}{@{}lcc@{}}
\toprule
\textbf{Model} & \textbf{Compiled} & \textbf{Base passed} \\
\midrule
Qwen2.5-Coder-1.5B-Instruct-GGUF    & 30 / 30 & 30 / 30 \\
DeepSeek-Coder-1.3B-Instruct-GGUF   & 30 / 30 & 30 / 30 \\
\bottomrule
\end{tabular}
\end{table}

\subsection{Branch Coverage}

Table~\ref{tab:cov} aggregates the per-task JaCoCo output from the CI
workflow. The improved suites close every avoidable coverage gap; the two
remaining Qwen-side misses sit on Java/001 and Java/006, where the
Qwen-generated parser contains a private helper-exit branch unreachable
through the public API under the dataset's contract.

\begin{table}[!t]
\renewcommand{\arraystretch}{1.15}
\caption{Aggregate branch coverage on the 30 HumanEval-X Java tasks.}
\label{tab:cov}
\centering
\footnotesize
\begin{tabular}{@{}lcccc@{}}
\toprule
& \multicolumn{2}{c}{\textbf{Base suite}} & \multicolumn{2}{c}{\textbf{Improved suite}} \\
\cmidrule(lr){2-3}\cmidrule(lr){4-5}
\textbf{Model}    & \textbf{Cov./Tot.} & \textbf{\%} & \textbf{Cov./Tot.} & \textbf{\%} \\
\midrule
Qwen 1.5B    & 123 / 128 & 96.09 & 126 / 128 &  98.44 \\
DeepSeek 1.3B & 128 / 130 & 98.46 & 130 / 130 & 100.00 \\
\bottomrule
\end{tabular}
\end{table}

\subsection{Mutation Score}

Because the mutants are encoded as JUnit assertions, the improved suites
trivially \emph{kill} every mutant we crafted and the base suites kill
\emph{none}. The derived mutation-kill rate per task is therefore $0/N$
for the LLM-produced base suites and $N/N$ for the improved suites, where
$N \in \{1,2\}$ is the number of mutants per task. Phase~2 supplements
this with PITest-driven automated mutation analysis on the composite
\texttt{BookScan} class.

\subsection{Model-Specific Behavioural Divergence}

Several tasks produced \emph{semantically distinct} implementations for
Qwen and DeepSeek (Table~\ref{tab:divergence}). The improved suites carry
different assertions for those tasks so that the suite documents what each
LLM actually did.

\begin{table}[!t]
\renewcommand{\arraystretch}{1.15}
\caption{Tasks where Qwen and DeepSeek diverge on observable behaviour.}
\label{tab:divergence}
\centering
\footnotesize
\begin{tabular}{@{}lp{6.0cm}@{}}
\toprule
\textbf{Task} & \textbf{Behavioural disagreement} \\
\midrule
Java/2  & \texttt{Math.floor} (Qwen) vs \texttt{(int)} cast (DeepSeek) for negative input. \\
Java/4  & \texttt{Optional.orElse(0.0)} (Qwen) vs division-by-zero \(\to\) \texttt{NaN} (DeepSeek) on empty list. \\
Java/19 & \texttt{indexOf == -1} (Qwen) vs \texttt{HashMap.get == null} \(\to\) NPE (DeepSeek). \\
Java/20 & First tied pair (Qwen) vs accumulate-all-ties (DeepSeek). \\
Java/21 & Returns a fresh list (Qwen) vs mutates input in place (DeepSeek). \\
\bottomrule
\end{tabular}
\end{table}

% =====================================================================
\section{Discussion}\label{sec:discuss}
% =====================================================================

\subsection{Comparison Between Models}

The two models reached comparable headline numbers. DeepSeek edges out
Qwen on aggregate branch coverage, but the gap is small enough that it
could plausibly be reproduced by a different sampling seed. The more
interesting differences are qualitative: as Table~\ref{tab:divergence}
shows, the two models pick different defensive coding styles for
negative-input edge cases, and those choices --- not the headline coverage
delta --- determine whether the generated code is robust under unusual
inputs.

\subsection{Effectiveness of the Test Strategies}

Three test strategies were applied in sequence: the dataset base test
(LLM-generated), the equivalence-class extension and the hand-crafted
mutation guards. Their effects are different in kind. The base test is
sufficient to catch obvious regressions but reaches its branch-coverage
ceiling quickly. The equivalence-class extension closes most of the
remaining gap, especially for tasks with non-trivial invalid-input or
empty-input behaviour. The hand-crafted mutation guards are the only
stage that detects \emph{behavioural} regressions that branch coverage
cannot --- e.g.\ replacing \verb|<| with \verb|<=| in a comparison, or
replacing \texttt{Math.floor} with an \texttt{(int)} cast. This is
consistent with the recent
literature~\cite{metatestgen,yang2024,lit_review}: coverage is necessary
but not sufficient.

\subsection{Branch Coverage versus Mutation Score}

The two metrics measure different things. Branch coverage tells us what
\emph{was executed}; the mutation guards tell us whether the suite would
\emph{notice a behavioural change}. A suite can reach 100\,\% branch
coverage and still kill 0\,\% of meaningful mutants, because the
assertions can be too lax. Conversely, a suite can kill all the
hand-crafted mutants and still leave a private helper unreached. We
therefore report both metrics, and Phase~2 will additionally use PITest to
give a third, automated mutation perspective.

\subsection{Threats to Validity}

\emph{Internal validity} is limited by local-model non-determinism; we
mitigate this by pinning the sampling seed and the model checkpoints, and
by re-running the entire workflow under GitHub Actions on every push.
\emph{External validity} is limited by the 30-prompt subset; HumanEval-X
samples algorithmic puzzles rather than full-fledged Java systems, so the
headline numbers should be read as ``what does this LLM look like on
bite-sized Java tasks'' rather than ``in production''. \emph{Construct
validity} depends on the coverage and mutation-kill proxies: coverage is
objective but blunt; the mutation set is hand-picked and is best read as
a \emph{lower bound} on the suite's robustness. Phase~2 supplements this
with PITest, which removes the human-curation bias from the mutation set.

% =====================================================================
\section{Literature Review}\label{sec:lit}
% =====================================================================

The literature set was selected to cover five complementary perspectives
on LLM-driven test generation, all published within the last three years
(2023--2026): an early end-to-end tool paper, an industrial deployment
study, a broad empirical evaluation of open-source LLMs, a hybrid
coverage-driven test-generation technique, and a recent systematic
review. The selection was performed by hand using Google Scholar and IEEE
Xplore; no LLM was used to surface or summarise the papers.

\textbf{ChatUniTest~\cite{chatunitest}} is an early reference for LLM-based
unit-test generation. Its main contribution is wrapping the model inside
a generation--validation--repair loop. The paper argues that raw
generations are often incomplete or syntactically broken, so validation
and repair are necessary rather than optional. This directly motivates
our normalisation plus repair loop in Phase~1.

\textbf{Meta TestGen-LLM~\cite{metatestgen}} shifts the problem from
generating tests from scratch to improving existing human-written tests.
Its lesson is methodological: model output is accepted only after it
passes build, execution, flakiness and coverage-gain filters. Our
pipeline uses the same instinct.

\textbf{Yang \emph{et al.}~\cite{yang2024}} is a direct empirical
comparison of multiple open-source LLMs on Java unit-test generation.
Two findings are relevant: prompt structure materially affects compilation
and coverage outcomes, and open-source models can be competitive but
still need additional engineering. Both observations are mirrored by our
own measurements.

\textbf{Panta~\cite{panta2024}} pushes beyond prompt tuning by combining
LLM generation with static control-flow analysis and dynamic coverage
feedback. The paper targets exactly the weakness that appears in our
project: coverage gaps on branch-heavy code. Our improved-test step is a
simpler instance of the same philosophy.

\textbf{Tasarsu, Tokmak and Catal~\cite{lit_review}} synthesise 38
peer-reviewed studies from 2020--2025 in a 2026 \emph{Cluster Computing}
systematic literature review. Their findings confirm patterns visible in
our project as well: coverage remains the dominant evaluation metric,
post-processing and integration logic are common, and
maintainability/robustness are still weaker than simple coverage
reporting. We complement coverage with hand-crafted mutation guards in
Phase~1 and with automated PITest in Phase~2 specifically because of this
observation.

% =====================================================================
\section{Conclusion}\label{sec:conclusion}
% =====================================================================

Phase~1 of the BLG~475E term project established a reproducible,
semi-agentic pipeline that runs two local code-oriented LLMs over 30
prompts from the HumanEval-X Java corpus. After normalization, JUnit~6
base-test conversion, equivalence-class-driven test improvement, and
hand-crafted mutation guards, both models reach $30/30$ base-test
success; the improved suites raise aggregate branch coverage from
$96.09\,\%$ / $98.46\,\%$ on the base suites to $98.44\,\%$ / $100.00\,\%$;
every (task, model) pair carries at least one \texttt{improvedMutation\ldots}
JUnit method. Seven tasks initially failed and were repaired through a
logged refactoring round.

The most useful methodological lesson from Phase~1 is that branch coverage
alone is insufficient: a suite can reach a high coverage percentage and
still fail to detect a meaningful operator-level mutation. A structured
combination of equivalence-class partitioning, boundary-value analysis and
mutation guards materially strengthens the resulting suite, and this
discipline carries forward into the Phase~2 integration testing of the
composite \texttt{BookScan} class.

% =====================================================================
\appendices
\section{Selected HumanEval-X Java Prompts}\label{app:prompts}
% =====================================================================

The 30 prompts selected for Phase~1, with their difficulty labels and
Phase~2 dependency flag, are listed in Table~\ref{tab:prompts-full}.

\begin{table}[!ht]
\renewcommand{\arraystretch}{1.05}
\caption{Phase 1 selected prompts (full list).}
\label{tab:prompts-full}
\centering
\footnotesize
\begin{tabular}{@{}llll@{}}
\toprule
\textbf{Task} & \textbf{Method} & \textbf{Diff.} & \textbf{P2 dep?} \\
\midrule
Java/2  & truncateNumber       & easy  & no \\
Java/3  & belowZero            & easy  & no \\
Java/9  & rollingMax           & easy  & no \\
Java/11 & stringXor            & easy  & no \\
Java/13 & greatestCommonDivisor& easy  & no \\
Java/14 & allPrefixes          & easy  & no \\
Java/22 & filterIntegers       & easy  & no \\
Java/23 & strlen               & easy  & \textbf{yes} \\
Java/24 & largestDivisor       & easy  & no \\
Java/27 & flipCase             & easy  & \textbf{yes} \\
\midrule
Java/0  & hasCloseElements     & moderate & no \\
Java/4  & meanAbsoluteDeviation& moderate & no \\
Java/5  & intersperse          & moderate & no \\
Java/8  & sumProduct           & moderate & no \\
Java/15 & stringSequence       & moderate & no \\
Java/16 & countDistinctChars   & moderate & no \\
Java/18 & howManyTimes         & moderate & \textbf{yes} \\
Java/19 & sortNumbers          & moderate & no \\
Java/25 & factorize            & moderate & no \\
Java/29 & filterByPrefix       & moderate & no \\
\midrule
Java/1  & separateParenGroups  & hard & no \\
Java/6  & parseNestedParens    & hard & no \\
Java/7  & filterBySubstring    & hard & no \\
Java/10 & isPalindrome         & hard & no \\
Java/12 & longest              & hard & no \\
Java/17 & parseMusic           & hard & no \\
Java/20 & findClosestElements  & hard & no \\
Java/21 & rescaleToUnit        & hard & no \\
Java/31 & isPrime              & hard & no \\
Java/39 & primeFib             & hard & no \\
\bottomrule
\end{tabular}
\end{table}

\section{Reproduction Recipe}\label{app:repro}

The Phase 1 pipeline is reproducible end-to-end:
\begin{lstlisting}[language=bash, caption=Reproducing the Phase 1 coverage CSVs.]
python Phase1/scripts/python/run_base_coverage.py \
  --tasks-root Phase1/tasks \
  --results   Phase1/results/base_coverage_results.csv \
  --java      $JAVA_HOME/bin/java \
  --javac     $JAVA_HOME/bin/javac \
  --junit-console  .tools/junit-platform-console-standalone-6.0.3.jar \
  --jacoco-agent   .tools/org.jacoco.agent-0.8.12-runtime.jar \
  --jacoco-cli     .tools/org.jacoco.cli-0.8.12-nodeps.jar
\end{lstlisting}

A second invocation with \verb|--test-stage improved| produces the
improved-coverage CSV. The same flow is automated in
\texttt{.github/workflows/phase1-ci.yml}.

% =====================================================================
\section*{Acknowledgment}
% =====================================================================

The two-person group split the Phase~1 workload along two complementary
tracks; both authors contributed equally and continuously cross-reviewed
each other's deliverables. \textbf{Mustafa Eren Ko\c{c}} (150190805) led the
\emph{engineering and automation} track: repository scaffolding, GitHub
Actions CI, the PowerShell + Python~3.11 driver layer, the local
\texttt{llama.cpp} execution pipeline, dataset-to-JUnit adapters, JaCoCo
automation and result CSV aggregation, the repair-loop driver, and the Git
commit discipline. \textbf{Onat Bar{\i}\c{s} Ercan} (150210075) led the
\emph{quality and analysis} track: 30-prompt selection rationale and
difficulty categorisation, per-task equivalence-partitioning and
boundary-value documents, test-smell audit, hand-crafted
mutation-operator catalogue and \texttt{improvedMutation\ldots} JUnit
method authoring, the five-paper literature review, and the IEEE journal
manuscript.

% =====================================================================
\begin{thebibliography}{9}
% =====================================================================

\bibitem{chatunitest}
Z.~Xie, Y.~Chen, C.~Zhi, S.~Deng, and J.~Yin, ``ChatUniTest: a
ChatGPT-based automated unit test generation tool,'' \emph{CoRR},
abs/2305.04764, 2023, doi: 10.48550/arXiv.2305.04764.

\bibitem{metatestgen}
N.~Alshahwan \emph{et al.}, ``Automated unit test improvement using large
language models at Meta,'' in \emph{Companion Proc.\ 32nd ACM Int.\ Conf.
Found.\ Softw.\ Eng. (FSE Companion)}, 2024, pp.~185--196, doi:
10.1145/3663529.3663839.

\bibitem{yang2024}
L.~Yang \emph{et al.}, ``On the evaluation of large language models in
unit test generation,'' in \emph{Proc.\ 39th IEEE/ACM Int.\ Conf.\
Automated Softw.\ Eng. (ASE)}, 2024, pp.~987--999, doi:
10.1145/3691620.3695529.

\bibitem{panta2024}
S.~Gu, N.~Nashid, and A.~Mesbah, ``LLM test generation via iterative
hybrid program analysis,'' \emph{CoRR}, abs/2503.13580, 2025. Accepted to
ICSE 2026, doi: 10.48550/arXiv.2503.13580.

\bibitem{lit_review}
M.~Tasarsu, A.~V.\ Tokmak, and C.~Catal, ``Test case generation using
large language models: a systematic literature review,'' \emph{Cluster
Computing}, vol.~29, no.~4, Springer, 2026, doi:
10.1007/s10586-026-06021-z.

\bibitem{humaneval}
M.~Chen \emph{et al.}, ``Evaluating large language models trained on
code,'' \emph{arXiv preprint arXiv:2107.03374}, 2021.

\bibitem{jnose}
J.~Aniche, B.~Alves, and J.~Fonseca, ``JNose: a JUnit test smell detection
tool,'' \emph{Online}: \url{https://jnosetest.github.io/}, accessed
2026-04-25.

\end{thebibliography}

\end{document}
